\documentclass[reqno,11pt]{amsart}
\usepackage{math327-lecture}
\usepackage{needspace}
\newcommand{\Cl}{\operatorname{Cl}}
\lectureinfo{23}{Nov. 2, 2026}{Permutation Classes and the Simplicity of the Alternating Group A5}

% Based on handwritten Lecture 20, Oct. 22, 2025, pp. 1--10,
% and Lecture 21, Oct. 24, 2025, pp. 1--11.
\begin{document}
\raggedbottom
\makelecturetitle

\section{Normal subgroups and conjugacy classes}

Recall that conjugation partitions a finite group into classes
$\Cl_G(x)=\{gxg^{-1}:g\in G\}$. Its stabilizer is the
centralizer $C_G(x)=\{g\in G:gx=xg\}$, and
\[
  |\Cl_G(x)|=[G:C_G(x)],\qquad
  |G|=\sum_{\text{distinct classes }C}|C|.
\]
Each class size divides $|G|$, and the identity contributes the
singleton class $\{1\}$.

\begin{proposition}\label{prop:normal-union}
A subgroup $N\leq G$ is normal if and only if it is a union of
conjugacy classes of $G$.
\end{proposition}
\begin{proof}
Normality says that $gng^{-1}\in N$ for all $n\in N$ and $g\in G$.
Equivalently, $\Cl_G(n)\subseteq N$ for every $n\in N$.
\end{proof}

For finite $G$, a normal subgroup's order is therefore a sum of
class sizes that includes the identity class. It must also divide
$|G|$ by Lagrange's theorem. These two requirements can be restrictive.

\begin{example}[Normal subgroups of $S_3$]
The class sizes in $S_3$ are $1,2,3$. A normal subgroup must include
the identity, so its possible orders from unions of classes are
$1,3,4,6$. Of these, only $1,3,6$ divide $6$.
Thus the only nontrivial proper normal subgroup is
\[
  A_3=\{1,(1\,2\,3),(1\,3\,2)\}.
\]
\end{example}

\section{Cycle decompositions and cycle types}

Every permutation is a product of disjoint cycles. Cycles are
\emph{disjoint} when they have no entries in common; disjoint cycles
commute. The decomposition is unique up to reordering the cycles
and cyclically shifting the entries within each cycle.
Fixed points may be recorded as cycles of length $1$.

\begin{example}
The cycles $(1\,2)$ and $(3\,4)$ are disjoint, whereas
$(1\,3)$ and $(3\,2)$ are not. A permutation can also be expressed
as a product of transpositions that need not be disjoint, for example
\[
  (1\,2\,3)=(2\,3)(1\,3).
\]
Products are composed from right to left.
\end{example}

\begin{lemma}[Parity and order of cycles]
A cycle of length $r$ has order $r$ and sign $(-1)^{r-1}$.
The order of a product of disjoint cycles of lengths
$r_1,\ldots,r_k$ is $\operatorname{lcm}(r_1,\ldots,r_k)$.
\end{lemma}
\begin{proof}
An $r$-cycle first returns every entry to its original position
after $r$ applications. It can be written as $r-1$ transpositions:
\[
  (i_1\,i_2\,\cdots\,i_r)
  =(i_1\,i_r)(i_1\,i_{r-1})\cdots(i_1\,i_2).
\]
For disjoint cycles, a power is the identity exactly when its
exponent is divisible by every cycle length.
\end{proof}

In particular, cycles of odd length are even permutations, and
cycles of even length are odd permutations.

\begin{theorem}[Conjugacy in $S_n$]\label{thm:cycle-type}
Two permutations in $S_n$ are conjugate if and only if they have
the same cycle type, including fixed points.
\end{theorem}
\begin{proof}
Conjugation relabels a cycle:
\[
  g(i_1\,\cdots\,i_r)g^{-1}=(g(i_1)\,\cdots\,g(i_r)).
\]
It therefore preserves the lengths of disjoint cycles.
Conversely, for two permutations of the same cycle type, match
their cycles of equal lengths and define $g$ by matching the
entries in those cycles in order. The displayed formula then
shows that $g$ conjugates one permutation to the other.
\end{proof}

\Needspace{18\baselineskip}
\begin{example}[The class equation of $S_4$]
The possible cycle types and their sizes are as follows.
\begin{center}
\begin{tabular}{ccl}
  \textbf{Representative} & \textbf{Class size} & \textbf{Count}\\ \hline
  $1$ & $1$ & Identity\\[2pt]
  $(1\,2)$ & $6$ & $\binom42$\\[2pt]
  $(1\,2\,3)$ & $8$ & $\binom43\cdot2$\\[2pt]
  $(1\,2\,3\,4)$ & $6$ & $3!$\\[2pt]
  $(1\,2)(3\,4)$ & $3$ & Three pairings of four entries
\end{tabular}
\end{center}
For instance, after choosing the entries of a $3$-cycle there are
two cyclic orders. For a $4$-cycle, putting $1$ first leaves $3!$
orders for the other entries. The class equation is
\[
  24=1+6+8+6+3.
\]
\end{example}

\Needspace{20\baselineskip}
\begin{example}[The class equation of $S_5$]\label{ex:s5}
There are seven cycle types in $S_5$.
\begin{center}
\begin{tabular}{cccl}
  \textbf{Representative} & \textbf{Class size} & \textbf{Parity} & \textbf{Count}\\ \hline
  $1$ & $1$ & Even & Identity\\[2pt]
  $(1\,2)$ & $10$ & Odd & $\binom52$\\[2pt]
  $(1\,2\,3)$ & $20$ & Even & $\binom53\cdot2$\\[2pt]
  $(1\,2\,3\,4)$ & $30$ & Odd & $\binom54\cdot3!$\\[2pt]
  $(1\,2\,3\,4\,5)$ & $24$ & Even & $4!$\\[2pt]
  $(1\,2)(3\,4)$ & $15$ & Even & $\binom54\cdot3$\\[2pt]
  $(1\,2\,3)(4\,5)$ & $20$ & Odd & $\binom53\cdot2$
\end{tabular}
\end{center}
Thus
\[
  120=1+10+20+30+24+15+20.
\]
The even permutations have total size $1+20+24+15=60$.
This is not yet the class equation of $A_5$: conjugation in $A_5$
allows only even conjugating permutations.
\end{example}

\section{Comparing classes in S\texorpdfstring{$_n$}{n} and A\texorpdfstring{$_n$}{n}}

For $\sigma\in A_n$, we have
\[
  \Cl_{A_n}(\sigma)\subseteq\Cl_{S_n}(\sigma),\qquad
  C_{A_n}(\sigma)=C_{S_n}(\sigma)\cap A_n.
\]
The first inclusion can be strict: a class in $S_n$ may split when
we restrict conjugation to $A_n$.

\begin{lemma}[Splitting of conjugacy classes]\label{lem:splitting}
Let $n\geq2$ and $\sigma\in A_n$.
\begin{enumerate}[label=\textbf{(\arabic*)}]
  \item If $C_{S_n}(\sigma)\subseteq A_n$, then
  $\Cl_{S_n}(\sigma)$ splits into two $A_n$-classes of equal size, and
  \[
    |\Cl_{A_n}(\sigma)|=\tfrac12|\Cl_{S_n}(\sigma)|.
  \]
  \item If $C_{S_n}(\sigma)$ contains an odd permutation, then
  $\Cl_{A_n}(\sigma)=\Cl_{S_n}(\sigma)$.
\end{enumerate}
\end{lemma}
\begin{proof}
Put $H=C_{S_n}(\sigma)$ and $K=H\cap A_n=C_{A_n}(\sigma)$.
Orbit--stabilizer gives
\[
  |\Cl_{S_n}(\sigma)|=\frac{n!}{|H|},\qquad
  |\Cl_{A_n}(\sigma)|=\frac{n!/2}{|K|}.
\]
If $H\subseteq A_n$, then $K=H$, giving the factor $1/2$.
Fix an odd permutation $t$. Every conjugator in $S_n$ is either
even or belongs to the coset $A_nt$, so the $S_n$-class is the union
of the $A_n$-classes of $\sigma$ and $t\sigma t^{-1}$.
Conjugation by $t$ gives a bijection between these two classes,
so they have equal size. The factor $1/2$ shows that they are distinct.

If $H$ contains an odd permutation, the sign map on $H$ is
surjective with kernel $K$. Hence $|H|=2|K|$, and the two class
sizes are equal. The inclusion then gives equality of the classes.
Equivalently, an odd conjugator can be multiplied by an odd
element of $H$ to obtain an even conjugator with the same effect
on $\sigma$.
\end{proof}

\Needspace{15\baselineskip}
\begin{example}[Two $5$-cycles that are not conjugate in $A_5$]
Let $c=(1\,2\,3\,4\,5)$ and $d=(1\,2\,3\,5\,4)$.
They are conjugate in $S_5$ because $d=(4\,5)c(4\,5)^{-1}$.
The centralizer of $c$ in $S_5$ is
\[
  C_{S_5}(c)=\langle c\rangle.
\]
Indeed, a permutation commuting with $c$ is determined by its
value at $1$, since $h(c^j(1))=c^j(h(1))$. There are at most five
choices, and the five powers of $c$ all commute with $c$.
These powers are even.

If an even $h$ also satisfied $hch^{-1}=d$, then
$(4\,5)^{-1}h$ would belong to $C_{S_5}(c)$, which is impossible
because it is odd. Thus $c$ and $d$ lie in distinct $A_5$-classes.
Each class has size $60/5=12$.
\end{example}

\section{The class equation and simplicity of A\texorpdfstring{$_5$}{5}}

The even cycle types in $S_5$ are the identity, $3$-cycles,
double transpositions, and $5$-cycles.
The centralizer of $(1\,2\,3)$ contains the odd permutation
$(4\,5)$, so its class does not split. Likewise, $(1\,2)$ is an
odd element centralizing $(1\,2)(3\,4)$, so that class does not
split. The $5$-cycles split into two classes as above.

\begin{example}[The class equation of $A_5$]
The classes have the following representatives and sizes.
\begin{center}
\begin{tabular}{ccc}
  \textbf{Representative} & \textbf{Class size} & \textbf{Centralizer order in $A_5$}\\ \hline
  $1$ & $1$ & $60$\\[2pt]
  $(1\,2\,3)$ & $20$ & $3$\\[2pt]
  $(1\,2)(3\,4)$ & $15$ & $4$\\[2pt]
  $(1\,2\,3\,4\,5)$ & $12$ & $5$\\[2pt]
  $(1\,2\,3\,5\,4)$ & $12$ & $5$
\end{tabular}
\end{center}
Hence
\[
  60=1+20+15+12+12.
\]
\end{example}

\begin{define}[Simple group]
A nontrivial group is \emph{simple} if its only normal subgroups
are $\{1\}$ and the whole group.
\end{define}

\begin{theorem}
The group $A_5$ is simple.
\end{theorem}
\begin{proof}
Let $N\trianglelefteq A_5$. By Proposition~\ref{prop:normal-union},
its order is $1$ plus a sum of some of $20,15,12,12$.
If $N\neq\{1\}$, then $|N|\geq13$. If it is also proper, the
only divisors of $60$ this large are $15,20,30$.
None of these is $1$ plus a sum of some of $20,15,12,12$.
Thus there is no nontrivial proper normal subgroup.
\end{proof}

\begin{remark}
More generally, $A_n$ is simple for every $n\geq5$.
The argument above proves the case $n=5$.
\end{remark}

\end{document}
